\begin{table}[t]
\centering
\caption{Geographic Sorting vs Gerrymandering: Representative States}
\label{tab:geographic_sorting}
\begin{tabular}{lrrrrr}
\toprule
\textbf{State} & \textbf{Enacted} & \textbf{Algorithmic} & \textbf{Premium} & \textbf{Geo Frac} & \textbf{Class} \\
& \textbf{EG (\%)} & \textbf{EG (\%)} & \textbf{(\%)} & \textbf{(\%)} & \\
\midrule
Maryland & -1.1 & +12.3 & -13.5 & 1795 & Geography Dominated \\
South Carolina & +0.5 & +5.4 & -4.9 & 1167 & Geography Dominated \\
Mississippi & +0.9 & -1.9 & +2.8 & 390 & Geography Dominated \\
New York & +25.8 & +13.5 & +12.4 & 58 & Mixed \\
Oregon & -5.4 & +1.6 & -7.0 & 57 & Mixed \\
Alabama & -6.0 & -2.6 & -3.3 & 41 & Mixed \\
Ohio & +12.8 & +3.3 & +9.5 & 30 & Gerrymandering Dominated \\
Iowa & -4.3 & +1.5 & -5.9 & 30 & Gerrymandering Dominated \\
Pennsylvania & +16.2 & -5.2 & +21.3 & 28 & Gerrymandering Dominated \\
North Carolina & +18.3 & -4.9 & +23.2 & 27 & Gerrymandering Dominated \\
\bottomrule
\end{tabular}
\begin{tablenotes}
\small
\item \textit{Enacted EG}: Efficiency gap in enacted districts. \textit{Algorithmic EG}: Efficiency gap in compact algorithmic districts (geographic baseline). \textit{Premium}: Gerrymandering premium (enacted - algorithmic). \textit{Geo Frac}: Average geographic fraction across three metrics (efficiency gap, mean-median, partisan bias). \textit{Class}: Geography-Dominated ($>$60\% geographic), Mixed (30-60\%), or Gerrymandering-Dominated ($<$30\% geographic).
\end{tablenotes}
\end{table}
